\section{Dependencias}

Para el correcto funcionamiento del solver, es necesario contar con los siguientes requisitos:

\begin{itemize}
    \item \textbf{Versión de Python}: Python 3.8 o superior (se recomienda Python 3.11+ por rendimiento). El framework es totalmente compatible con cualquier versión moderna de Python (incluyendo Python 3.12 y superiores) sin restricciones de compatibilidad.
    \item \textbf{IDEs recomendados}: Visual Studio Code, PyCharm.
    \item \textbf{Gestión de dependencias}: Se utiliza el módulo \texttt{venv}, incluido en la biblioteca estándar de Python (desde la versión 3.3), para crear entornos virtuales de forma aislada.
\end{itemize}

\textbf{Nota}: La creación de un entorno virtual es opcional. Si no desea aislar las dependencias del proyecto y prefiere instalar los paquetes de Python directamente en el entorno global del sistema, puede omitir los pasos de creación y activación del entorno virtual y proceder directamente a la sección \ref{sec:instalar_dependencias}. Sin embargo, se recomienda encarecidamente el uso de un entorno virtual para evitar conflictos entre versiones de librerías en diferentes proyectos.

\subsection{Crear y Activar el entorno virtual con venv}

El módulo \texttt{venv} permite crear entornos virtuales aislados para gestionar paquetes sin afectar otras configuraciones del sistema. Para ello:

1.  Navegue hasta el directorio de su proyecto en la terminal o línea de comandos.
2.  Cree un entorno virtual llamado \texttt{env} (puede elegir otro nombre si lo prefiere) con el siguiente comando:

    \begin{verbatim}
    python -m venv env 
    \end{verbatim}
    % O en algunos sistemas podría ser necesario usar python3:
    % python3 -m venv env

3.  Una vez creado, active el entorno virtual. El comando varía según su sistema operativo y terminal:
    \begin{itemize}
        \item \textbf{Windows (cmd.exe)}:
        \begin{verbatim}
env\Scripts\activate.bat
        \end{verbatim}
        \item \textbf{Windows (PowerShell)}:
        \begin{verbatim}
.\env\Scripts\Activate.ps1 
        \end{verbatim}
        (Nota: Es posible que necesite ajustar la política de ejecución de scripts en PowerShell con \texttt{Set-ExecutionPolicy -ExecutionPolicy RemoteSigned -Scope Process}).
        \item \textbf{Unix o MacOS (bash/zsh)}:
        \begin{verbatim}
source env/bin/activate
        \end{verbatim}
    \end{itemize}

Al activarse correctamente, la línea de comandos debería mostrar el nombre del entorno virtual entre paréntesis al inicio, similar a:

\begin{verbatim}
(env) C:\Users\Usuario\Proyecto>
\end{verbatim}
o
\begin{verbatim}
(env) usuario@computador:~/Proyecto$ 
\end{verbatim}


\subsection{Configurar Visual Studio Code}

Para asegurarse de que Visual Studio Code use el intérprete de Python correcto dentro del entorno virtual:

1.  Abra la \textbf{Command Palette} (Paleta de Comandos) con:
    \begin{verbatim}
    Ctrl + Shift + P 
    \end{verbatim}
    (o \texttt{Cmd + Shift + P} en macOS).
2.  Escriba \texttt{>Python: Select Interpreter} (o \texttt{>Python: Seleccionar Intérprete}) y presione Enter.
3.  Seleccione la opción que apunta al ejecutable de Python dentro de su entorno virtual. Usualmente contendrá la ruta \texttt{./env/bin/python} (Unix/macOS) o \texttt{.\env\Scripts\python.exe} (Windows). VS Code a menudo detecta y sugiere automáticamente los entornos \texttt{venv}.

\subsection{Instalar dependencias}
\label{sec:instalar_dependencias}

Para instalar las dependencias del proyecto, asegúrese de que el entorno virtual (\texttt{env}) esté activado (verá \texttt{(env)} al inicio de la línea de comandos) y ejecute en la consola:

\begin{verbatim}
pip install -r requirements.txt
\end{verbatim}